% ============================================================
%  Chapter -- Related Work
% ============================================================
\chapter{Related Work}
\label{ch:related}

This chapter surveys the body of prior work most relevant to the three
pillars of this thesis: learned point cloud compression, post-processing
of compressed 3D geometry, and rate-conditioned neural networks.

% ============================================================
\section{Learned Point Cloud Compression}
\label{sec:related-compression}

\subsection{Voxel-Based Learned Codecs}

The first generation of learned point cloud codecs, pioneered by
Quach et al.~\cite{quach2019learning} and Wang et al.~\cite{wang2021multiscale},
applied convolutional autoencoders to binary occupancy grids.
PCGCv2 introduced a multi-scale architecture:
the occupancy grid is processed by a cascade of sparse convolutional layers
at progressively coarser resolutions, each level encoding a latent that
is entropy-coded and transmitted.
The decoder reconstructs the occupancy grid from the bottom up, conditioned
on the decoded coarse representation.
This coarse-to-fine strategy allows the codec to allocate bits efficiently
across resolutions, with coarse geometry coded at low cost and fine details
refined at the finer levels.

SparsePCGC~\cite{wang2023sparsepcgc} extended PCGCv2 with fully sparse
tensor representations throughout the latent hierarchy, replacing dense
tensor operations with sparse ones at every scale.
This reduces memory and computation for large-scale point clouds and improves
rate-distortion efficiency.

Unicorn~\cite{wang2025unicorn} represents the current state of the art among
occupancy-based codecs, introducing universal multi-scale conditional coding
with cross-scale and cross-position entropy model conditioning.
Unicorn achieves strong rate-distortion performance at extreme compression
rates (bpp $< 0.02$) that earlier codecs struggle to reach.
Both SparsePCGC and Unicorn are relevant to this thesis: SparsePCGC is the
direct architectural successor to PCGCv2 and shares its oversmoothing
characteristics; Unicorn serves as the unseen-codec target for zero-shot
transfer experiments (Section~\ref{sec:generalisation}).

\subsection{Density-Adaptive and Point-Based Methods}

An alternative to voxel-based approaches is to operate directly on
continuous point coordinates.
He et al.~\cite{he2022density} identified the density-aware bottleneck:
existing codecs fail to adapt their bit allocation to spatially variable
point densities, leading to quality imbalances similar to the oversmoothing
documented in this thesis (though they focus on density rather than curvature).
Their density-adaptive codec explicitly models density variation in the
entropy model.

Point-cloud codecs based on learned local feature descriptors (folding
networks~\cite{yang2018foldingnet}, flow-based
models~\cite{yang2019pointflow}) have been proposed but generally lag behind
occupancy-based methods on the MPEG evaluation benchmarks.

\subsection{MPEG Standardisation Activity}

The MPEG standardisation process has evaluated both G-PCC~\cite{graziosi2020gpcc}
and V-PCC~\cite{schwarz2018emerging} as the two principal traditional
standards.
G-PCC covers geometry compression for sparsely sampled (LiDAR) content
via octree coding, while V-PCC addresses dense volumetric video via 2D atlas
projection.
The call for proposals on learned point cloud compression~\cite{mpeg2024cfp}
(MPEG-PCC CfP, 2024) evaluated several learned codecs including SparsePCGC
and Unicorn against the G-PCC baseline, with learned codecs showing
consistent advantages at low bitrates.

% ============================================================
\section{Post-Processing for Compressed 3D Content}
\label{sec:related-postprocessing}

Table~\ref{tab:related-comparison} positions the approach of this thesis
relative to the most closely related post-processing methods along the
key design dimensions.

\begin{table}[t]
  \centering
  \caption{%
    Comparison of learned post-processing approaches for compressed 3D geometry.
    \textbf{Standalone}: no access to codec internals required.
    \textbf{Multi-rate}: single model covers multiple rate points.
    \textbf{Learned codec}: targets a learned (occupancy) codec.
  }
  \label{tab:related-comparison}
  \small
  \begin{tabular}{lcccl}
    \toprule
    \textbf{Method} & \textbf{Standalone} & \textbf{Multi-rate} & \textbf{Learned codec} & \textbf{Artifact type} \\
    \midrule
    Fan et al.~\cite{fan2022deep}     & \checkmark & $\times$ & $\times$ & G-PCC quant. noise \\
    GRASP-Net~\cite{pang2022graspnet} & \checkmark & $\times$ & $\times$ & G-PCC quant. noise \\
    Guarda et al.~\cite{guarda2025deep} & $\times$ & \checkmark & \checkmark & Learned codec \\
    Xu et al.~\cite{xu2024fast}       & $\times$ & \checkmark & \checkmark & Learned codec \\
    \textbf{Ours (FiLM post-proc.)}   & \checkmark & \checkmark & \checkmark & Oversmoothing \\
    \bottomrule
  \end{tabular}
\end{table}

\subsection{Post-Processing for G-PCC Geometry}

Fan et al.~\cite{fan2022deep} proposed a graph neural network for
post-processing G-PCC decoded geometry.
The network constructs a local graph on the decoded point cloud and applies
graph convolutions to estimate per-point corrections.
On the MPEG evaluation sequences, the method achieves BD-PSNR improvements
of 1--2~dB.
However, this approach is designed for G-PCC's quantisation-noise artifact
profile and does not address the spatially biased oversmoothing of learned
codecs.

Pang et al.~\cite{pang2022graspnet} proposed GRASP-Net, a residual analysis
and synthesis network for G-PCC geometry enhancement.
GRASP-Net uses a graph-based local feature extractor and a residual decoder
to predict point-level corrections from the decoded cloud.
Like Fan et al., this is tailored to G-PCC artifacts.

\subsection{Integrated Refinement in Learned Codecs}

Guarda et al.~\cite{guarda2025deep} and Xu et al.~\cite{xu2024fast} propose
refinement modules that are integrated directly into the learned codec
decoder, rather than operating as a standalone post-processor.
These approaches modify the codec architecture and require joint training of
the codec and refinement components.
They achieve stronger per-codec improvements but cannot be applied to
pre-trained or black-box codecs.
Our approach, by contrast, is a standalone post-processor that requires no
access to the codec internals and can be applied to any codec's output.

\subsection{Attribute Enhancement}

Concurrent work on attribute (colour) post-processing~\cite{ugae2025} targets
colour restoration after geometric compression.
The method takes the decoded point cloud (with colour attributes) and a
reference geometry (if available) to enhance colour fidelity.
This is complementary to geometric post-processing: an ideal system would
jointly refine both geometry and colour.

\subsection{Analogy with Image Compression Artefact Removal}

The problem of removing compression artifacts from images has a longer
history.
Deep learning approaches (ARCNN~\cite{dong2015compression}, DnCNN~\cite{zhang2017dncnn})
demonstrated that convolutional networks can effectively reduce JPEG blocking
and ringing artifacts, learning to invert the characteristic distortions of
DCT-based compression.
The analogy motivates the learned post-processing approach for point cloud
codecs, but the 3D sparse structure and the non-uniform spatial distribution
of artifacts (concentrated at edges, not uniformly distributed) require
specialised architectures.

% ============================================================
\section{Rate-Conditioned and Multi-Rate Neural Networks}
\label{sec:related-rate}

\subsection{Variable-Rate Learned Image Compression}

In learned image compression, achieving multiple rate-distortion operating
points with a single model (rather than training separate models per rate)
is an active research area.
Choi et al.~\cite{choi2019variable} proposed conditioning the entropy model
on a quantisation step size to achieve variable-rate compression.
Cui et al.~\cite{cui2021asymmetric} introduced asymmetric gained autoencoders,
using learned gain factors to scale the latent representation.
Song et al.~\cite{song2021variable} and Li et al.~\cite{li2024neural} further
developed variable-rate mechanisms based on gain units and neural substitutions.

In each case, a scalar or vector conditioning signal shifts the codec's
rate-distortion operating point without retraining.
FiLM conditioning serves the same purpose here: a scalar rate input modulates
intermediate features via affine transformation, and a single model covers a
range of operating points.

\subsection{Feature-wise Linear Modulation in Context}

FiLM~\cite{perez2018film} was originally introduced for visual question
answering, where language embeddings modulate visual features to enable
compositional reasoning.
The closest structural precedent to the present work is multi-domain
adaptation~\cite{rebuffi2018efficient}, where a shared backbone handles
multiple visual domains by inserting per-domain FiLM layers --- the same
parameter-sharing argument that motivates using a single model across
rate-distortion operating points here.
TADAM~\cite{oreshkin2018tadam} used the same mechanism for few-shot learning:
a task embedding predicted at meta-test time steers intermediate features
without updating backbone weights, demonstrating that FiLM provides a
clean interface between a discrete or continuous conditioning variable and
a fixed-backbone computation.
In this thesis the conditioning signal is the log-BPP scalar of the decoded
cloud, and the goal is for a single model to cover the full range of
compression rates without per-rate weight copies.

\subsection{Multi-Task Learning and Loss Weighting}

Training a model on multiple related tasks or multiple operating conditions
(as in multi-rate training) raises the question of how to balance task losses.
Kendall et al.~\cite{kendall2018multi} proposed weighting losses by the
homoscedastic uncertainty of each task, treating the task weights as learned
parameters.
Chen et al.~\cite{chen2018gradnorm} proposed GradNorm, normalising task
gradient magnitudes to ensure balanced learning rates.
These methods assume that tasks can decouple their learning dynamics,
which requires either task-specific parameters or batch-level gradient
separation.

In this work, all rates share the backbone, preventing decoupled learning
dynamics.
The ablation in Section~\ref{ssec:ablation-weighting} confirms that adaptive
weighting schemes are not beneficial in this setting, and that the correct
prior --- expressed by static weights $[1.0, 0.3, 0.05]$ --- outperforms
all alternatives.

% ============================================================
\section{Quality Metrics for Point Cloud Compression}
\label{sec:related-metrics}

\subsection{Standard MPEG Metrics}

Point cloud quality has been standardly measured by D1 PSNR
(point-to-point nearest-neighbour error)~\cite{tian2017geometric}
and D2 PSNR (point-to-plane error).
Both metrics aggregate error uniformly over the surface, without regard to
local geometric complexity.

\subsection{Perceptual Quality Studies}

Subjective quality studies~\cite{ak2024basics,prazeres2024quality} have
documented that standard aggregate metrics fail to predict human perceptual
quality judgements for learned codec output.
In particular, geometric artifacts that concentrate at high-curvature
regions (edges, ridges) are perceptually more annoying than equivalent-MSE
errors distributed uniformly.
These studies motivate the development of geometry-aware metrics that weight
errors by their perceptual relevance.

The curvature-stratified PSNR introduced in this thesis is a step in this
direction: it separates quality evaluation by surface geometry, exposing the
edge degradation that aggregate metrics mask.
A fully perceptual metric would require psychophysical experiments with human
observers; the curvature-stratified metric is a computationally tractable
approximation that correlates with perceptual importance.

\subsection{Graph-Based and Feature-Based Metrics}

Other works have proposed metrics based on local shape descriptors, point
normal consistency, and graph spectral
analysis~\cite{meynet2020pcqm,alexiou2018angular,yang2022graphsim}.
PCQM~\cite{meynet2020pcqm} computes a multi-scale feature distance that
correlates well with subjective quality judgements but requires building
local geometric descriptors and performing pairwise distance comparisons
at several spatial scales, with computation cost substantially exceeding
that of a single nearest-neighbour search.
GraphSIM~\cite{yang2022graphsim} constructs a graph over the decoded cloud
and compares local gradient distributions via spectral analysis, adding
graph construction and eigendecomposition overhead.
Neither metric has been integrated into the MPEG evaluation toolchain,
so their use requires custom implementations that are not reproducible
across evaluation sites and whose results are not comparable to the
published D1/D2 baselines.
The curvature-stratified PSNR introduced in this thesis avoids these
drawbacks: it extends the D1 nearest-neighbour computation---already
executed as part of any standard quality evaluation---with a single $k$-NN
PCA pass to assign each point to a curvature stratum.
The additional cost is proportional to one extra $k$-NN query, the
output is a pair of standard PSNR values that slot directly into existing
reporting tables, and the stratification directly targets the spatially
non-uniform artifact distribution that motivates this work.
